% TOP_PROBLEM: {"id": 2303001, "title": "Research Problems in Function Theory — Problem 3.1", "category_id": 8, "category": {"id": 8, "name": "analysis", "display_name": "Analysis", "description": "Limits, continuity, calculus, and function theory.", "slug": "analysis", "order_index": 8, "created_at": "2026-07-31T15:26:25.671Z"}, "status": "open"}
% FIRST_POSTED: null
% ENTRY_KIND: statement_only
% Imported from: batch12.tex; UnsolvedMath ID 2303001
% Compile twice: lualatex 2303001.tex
\documentclass[11pt]{article}
\usepackage{fontspec}
\setmainfont{Latin Modern Roman}
\usepackage{amsmath,amssymb,mathtools,unicode-math}
\setmathfont{Latin Modern Math}
\usepackage{xurl,hyperref}
\hypersetup{colorlinks=true,urlcolor=blue,pdftitle={UnsolvedMath Problem 2303001}}
\newcommand{\sref}[2]{\hyperlink{source:#1:#2}{[S#2]}}
\newcommand{\cataloguescope}[1]{\paragraph{Mathematical question.}#1}
\begin{document}
% UnsolvedMath ID 2303001; source catalogue code AMR-022-3001
\setcounter{section}{2303000}
\section{Research Problems in Function Theory — Problem 3.1}\label{2303001}
{\small\sffamily UnsolvedMath ID 2303001 · Analysis}
\subsection{Definitions and mathematical statement}

\paragraph{Shared notation.}$\mathbb N=\{1,2,\ldots\}$, and $\mathbb Z,\mathbb Q,\mathbb R,\mathbb C$ denote the integers, rationals, reals and complex numbers. Any other conventions are specified locally.


A subharmonic function is an upper semicontinuous function with values in $[-\infty,\infty)$, not identically $-\infty$ on a connected component, that satisfies the mean-value inequality on every closed ball in its domain. A real harmonic function is a twice continuously differentiable function $u$ satisfying $\Delta u=0$. A harmonic polynomial on $\mathbb R^2\cong\mathbb C$ is a polynomial in the two real coordinates that is harmonic. A path to infinity is a continuous map $\gamma:[0,\infty)\to\mathbb C$ with $|\gamma(t)|\to\infty$.
The source update gives counterexamples of Barth, Brannan and Hayman even for growth powers $\alpha\geq1/2$.
\paragraph{Statement recorded in the supplied source.}
For every harmonic function $u$ on the plane that is not a polynomial and every integer $n\geq1$, does there exist a path $\gamma_n$ to infinity with
\[\frac{u(\gamma_n(t))}{|\gamma_n(t)|^n}\longrightarrow+\infty?\]
Can one choose a single path $\gamma$ for which this holds for every integer $n\geq1$? \sref{2303001}{2}
\subsection{Short English statement}
Must a nonpolynomial harmonic function grow faster than any prescribed power along some path, and can one path work for all powers?
\subsection{Sources}
\begin{description}
\item[\hypertarget{source:2303001:1}{S1}] ULAM.AI and the UnsolvedMath Contributors, UnsolvedMath: A Curated Collection of Open Mathematics Problems, record 2303001 (AMR-022-3001). CC BY 4.0. \url{https://huggingface.co/datasets/ulamai/UnsolvedMath}.
\item[\hypertarget{source:2303001:2}{S2}] W. K. Hayman and E. F. Lingham, \emph{Research Problems in Function Theory}, arXiv:1809.07200v2 (2018), Problem 3.1 and Update 3.1. \url{https://arxiv.org/abs/1809.07200}.
\end{description}
\label{end:2303001}
\end{document}
